% Build: latexmk -lualatex -interaction=nonstopmode -halt-on-error global-stafford-dossier.tex
% Standalone companion, not a replacement manuscript. Source graph:
% global-stafford/docs/proof-graph.yaml at b5654dd; paper pin in PLAN.md.
\documentclass[11pt,a4paper]{article}
\usepackage[margin=2.2cm]{geometry}
\usepackage{fontspec}
\setmainfont{DejaVu Serif}
\setsansfont{DejaVu Sans}
\setmonofont{DejaVu Sans Mono}[Scale=.80]
\usepackage{amsmath,amssymb,amsthm,mathtools}
\usepackage{microtype,booktabs,array,longtable,fvextra}
\usepackage{tikz,pdflscape,xcolor}
\usetikzlibrary{arrows.meta,calc}
\definecolor{ink}{HTML}{203B55}
\definecolor{bluewash}{HTML}{E8F0F7}
\definecolor{amberwash}{HTML}{FFF0D1}
\PassOptionsToPackage{obeyspaces}{url}
\usepackage[colorlinks=true,linkcolor=ink,urlcolor=ink,bookmarksnumbered=true]{hyperref}
\hypersetup{pdftitle={Global Stafford: proof architecture and Challenge dossier},pdfauthor={Christopher Albert}}
\usepackage{xurl}
\newcommand{\code}[1]{\nolinkurl{#1}}
\newcommand{\D}{\mathcal D}
\newcommand{\ord}{\operatorname{ord}}
\newcommand{\backmap}{\hfill\hyperref[fig:map]{\small Return to proof map}}
\newcommand{\leanref}[2]{\par\smallskip\noindent\textbf{Lean.} \href{https://github.com/itpplasma/global-stafford-formal/blob/b21883a5b3d8f46922713049c3b060523ea3a771/#1}{\nolinkurl{#1}}.\par\noindent\small\code{#2}\par\normalsize\backmap}
\newtheorem*{theorem}{Theorem (Global Stafford)}
\setlength{\emergencystretch}{2em}
\setlength{\parskip}{.35em}
\begin{document}
\begin{landscape}
\thispagestyle{empty}
\phantomsection\label{fig:map}
{\LARGE\bfseries Global Stafford: the proof map}\hfill 8 September 2026
\par\smallskip
\noindent Christopher Albert\quad\textcolor{ink}{\textbullet}\quad Proof architecture and Challenge dossier
\par\smallskip
\noindent\small Every box links to its statement and Lean evidence below. Solid arrows show the main dependencies;
 the dashed arrow supplies a theorem hypothesis. Transitive infrastructure edges are omitted.
\begin{center}
% Projection of global-stafford/docs/proof-graph.yaml at b5654dd.
% Node claims and dependencies follow source_descent_candidate and the
% gs-* nodes. Conjugation uses the recorded Remark 4.2 formal_variant.
% Transitive infrastructure edges are omitted; dashed edge supplies the
% whole-chart hypothesis of the bounded-order source theorem.
\begin{tikzpicture}[x=1cm,y=1cm,
  claim/.style={draw=ink,fill=bluewash,line width=.9pt,rounded corners=2pt,
    text width=5.1cm,minimum height=1.1cm,align=center,font=\small},
  imported/.style={claim,fill=amberwash,double,double distance=.6pt},
  flow/.style={-{Stealth[length=2mm]},draw=ink,line width=.7pt},
  edgeword/.style={font=\scriptsize,fill=white,inner sep=2pt,align=center}]
\node[imported] (weyl) at (0,0) {\hyperref[sec:weyl]{\textbf{Weyl S38 over every field}\\IMPORTED / LEAN}};
\node[claim] (structure) at (8,0) {\hyperref[sec:structure]{\textbf{Smooth\\operator structure}\\LEAN · Phase II producers}};
\node[claim] (rank) at (16,0) {\hyperref[sec:rank]{\textbf{Finite-right-rank ascent}\\LEAN · Lemma 6.1}};
\node[claim] (order) at (0,-2.5) {\hyperref[sec:order]{\textbf{Ordered power clearance}\\LEAN · (1.1), (1.2), Lemma 1.1}};
\node[claim] (square) at (8,-2.5) {\hyperref[sec:square]{\textbf{Squared-annihilator extraction}\\LEAN · Section 3}};
\node[claim] (chart) at (16,-2.5) {\hyperref[sec:chart]{\textbf{Fixed chart over central $k(t)$}\\LEAN · Lemma 6.2}};
\node[claim] (rho) at (0,-5.0) {\hyperref[sec:rho]{\textbf{Right-moved\\conjugation}\\LEAN · Remark 4.2}};
\node[claim] (bounded) at (8,-5.0) {\hyperref[sec:bounded]{\textbf{Bounded-order source family}\\LEAN · Theorem 4.1}};
\node[claim] (protect) at (0,-7.5) {\hyperref[sec:protect]{\textbf{Preservation\\of old charts}\\LEAN · Lemma 1.2}};
\node[claim] (patch) at (8,-7.5) {\hyperref[sec:patch]{\textbf{Finite-cover source patching}\\LEAN · Theorem 5.1}};
\node[claim,line width=1.4pt] (terminal) at (8,-10.0) {\hyperref[sec:terminal]{\textbf{Global Stafford: $1=dR+FdS$}\\LEAN · Theorem 6.3}};
\draw[flow] (structure) -- node[edgeword,sloped,above]{order filtration} (order);
\draw[flow] (structure) -- node[edgeword,right]{right Ore} (square);
\draw[flow] (structure) -- node[edgeword,sloped,above]{\'{e}tale coordinates} (chart);
\draw[flow] (rank) -- node[edgeword,right]{finite rank} (chart);
\draw[flow] (weyl.north) -- ++(0,.65) -| node[edgeword,pos=.25,above]{arbitrary field, including $k(t)$} (chart.north west);
\draw[flow] (order) -- node[edgeword,left]{binomial formulas} (rho);
\draw[flow] (square) -- node[edgeword,right]{written cofactors} (bounded);
\draw[flow] (rho) -- node[edgeword,above]{fixed polynomial\\$E'(t)$} (bounded);
\draw[flow,dashed] (chart) |- node[edgeword,pos=.72,above]{whole-ring $\mathrm{S38}_{\rm poly}$} (bounded);
\draw[flow] (order.west) -- ++(-.7,0) |- node[edgeword,pos=.65,left]{contractive\\error powers} (protect.west);
\draw[flow] (bounded) -- node[edgeword,right]{choose $N$ after $M,\ell,H$} (patch);
\draw[flow] (protect) -- node[edgeword,above]{preserve\\certificates} (patch);
\draw[flow] (patch) -- node[edgeword,right]{right clearance, B\'ezout} (terminal);
\draw[flow] (chart.east) -- ++(.65,0) |- node[edgeword,pos=.8,below]{fixed principal cover} (terminal.east);
\end{tikzpicture}

\par\smallskip\hyperref[sec:literature]{Literature and source roles for this map}\par
\end{center}
\noindent\small Blue solid boxes: project Lean proofs. Amber double box: imported Lean theorem.
These tags refer to the recorded kernel verification, not human mathematical review.
The map follows the research graph at \code{b5654dd}, using its Remark 4.2 formal variant.
\end{landscape}

\section{The theorem and the meaning of its statement}\label{sec:statement}
This formalization presents the Global Stafford proof developed in this
research project. Its informal exposition and internal working notes record
project chronology, including informal development before parts of the Lean
implementation. That chronology does not establish external historical priority.
The dated literature assessment is recorded in
\code{docs/provenance-literature-novelty.md}; references to the ``paper'' mean
the internal exposition.

This companion states the mathematical claim in TeX, compares it with the
Mathlib-only Challenge, and explains the mechanisms producing the global
certificate. Its correspondence uses the internal note
\code{notes/source-changing-descent-2026-09-07.md}. The informal argument
preceded parts of the Lean implementation; the dated record is retained below.

\begin{theorem}
Let $k$ be any field of characteristic zero and let $A$ be a smooth integral
finitely generated commutative $k$-algebra. For every nonzero
$d\in\D_k(A)$ there exist $F,R,S\in\D_k(A)$ such that
\[
                       1=dR+FdS.
\]
Here $\D_k(A)$ is the algebra of finite-order $k$-linear differential operators
on $A$, and multiplication is composition in the written order.
\end{theorem}

Write $m_a(x)=ax$ for multiplication by $a\in A$. For $P\in\operatorname{End}_k(A)$,
put $[P,a]=P\circ m_a-m_a\circ P$. Define
\[
\begin{aligned}
\D^{\le0}_k(A)&=\{P:[P,a]=0\text{ for every }a\in A\},\\
\D^{\le r+1}_k(A)&=\{P:[P,a]\in\D^{\le r}_k(A)\text{ for every }a\in A\},\\
\D_k(A)&=\bigcup_{r\ge0}\D^{\le r}_k(A).
\end{aligned}
\]
An order-zero operator is multiplication by $P(1)$, since commuting with
$m_a$ gives $P(a)=aP(1)$. This is Grothendieck's finite-order definition.
All operator ideals and modules in the argument are right-sided. Thus
$d\D+Fd\D=\D$ means precisely the displayed certificate; $R$ and $S$ are
right cofactors. The first generator remains the original $d$, and the
second generator is the left multiple $Fd$ of that same $d$.

There is no assumption that $k$ is algebraically closed, no positive-dimensional
restriction, and no global coordinate system on $\operatorname{Spec} A$.
For a polynomial algebra the statement specializes to the imported Weyl result.
The proof does not prescribe a Bernstein-degree bound on $F$ in the global case.

\paragraph{Exact source and verification scope.}
The internal-note revision is \nolinkurl{9a096dc9f6c990bdfb3e8a24110989648c70b3f3}, with
SHA-256 \nolinkurl{22f047887464d268c1def305e697232826799bc112cec812e31e2cb8f8172d03}.
It is unchanged at research revision \code{b5654dd}. The independently replayed
formal proof is \nolinkurl{b21883a5b3d8f46922713049c3b060523ea3a771}.
Release \code{v1.0.3} uses unique submission module names. Phase I/full replay and both standard and canonical-Challenge comparisons passed with the dependency pins below.
Phase I and II are complete in that record. Human expert review of the paper and
its statement correspondence remains open.

\paragraph{Literature assessment and registry record.}
The project's 9 September 2026 literature audit located no prior proof of the
exact same-divisor theorem. It records Bellamy's 2026 survey as still listing
the weaker smooth-affine two-generator question as open, and distinguishes
the three-generator result of Coutinho--Holland from the present statement.
These findings provide evidence of novelty; they do not certify absolute
historical priority. The audit gives the sources, search scope and limitations.
The repository links Palomar record
\href{https://palomar-registry.org/entry?id=PALOMAR-2026-09-05-000007&version=2}{\code{PALOMAR-2026-09-05-000007}, version 2}.
Registry status, machine checking and literature priority are distinct records.

\section{TeX statement compared with GlobalStaffordChallenge.lean}\label{sec:comparison}
The Challenge states the theorem on endomorphisms and explicitly requires
finite order for each witness. The substantive proof instead uses the
subalgebra of finite-order operators. These are two presentations of the same
objects, connected by proved equivalences.

\begin{longtable}{@{}>{\raggedright\arraybackslash}p{.25\textwidth}>{\raggedright\arraybackslash}p{.38\textwidth}>{\raggedright\arraybackslash}p{.31\textwidth}@{}}
\toprule
Mathematical statement & Challenge expression & Correspondence \\ \addlinespace[.4em]
\midrule\endhead
Every characteristic-zero field $k$ & \code{(k : Type u)}\newline \code{[Field k]} \code{[CharZero k]} & Arbitrary field in an arbitrary common universe; no closure assumption. \\ \addlinespace[.4em]
Smooth integral affine $A/k$ & \code{[CommRing A]} \code{[IsDomain A]}\newline \code{[Algebra k A]}\newline \code{[Algebra.Smooth k A]} & Smooth means formally smooth and finitely presented in the pinned Mathlib. Over a field this supplies finite type. \\ \addlinespace[.4em]
$P\in\operatorname{End}_k(A)$ & \code{P : Module.End k A} & $k$-linear maps $A\to A$; multiplication is composition. \\ \addlinespace[.4em]
$[P,a]=Pm_a-m_aP$ & \code{commutator P a} & Exactly the difference of those two ordered products. \\ \addlinespace[.4em]
$\ord P\le n$ & \code{IsOrderLE n P} & The recursive commutator definition above, including its order-zero clause. \\ \addlinespace[.4em]
$d\in\D_k(A)$, $d\ne0$ & \code{IsDifferentialOperator d → d ≠ 0 →} & Both hypotheses are explicit; $d=0$ is excluded. \\ \addlinespace[.4em]
$F,R,S\in\D_k(A)$ & \code{∃ F R S : Module.End k A,} followed by three \code{IsDifferentialOperator} predicates & No witness is allowed to be an infinite-order endomorphism. \\ \addlinespace[.4em]
$1=dR+FdS$ & \code{(1 : Module.End k A)}\newline \code{= d * R + F * d * S} & $1$ is the identity map. Associativity changes parentheses only, not factor order. \\ \addlinespace[.4em]
\bottomrule
\end{longtable}

\paragraph{The transport proof.}
\code{GlobalStafford.Assembly.isOrderLE_iff_mem_order} proves by induction on $n$
that the Challenge predicate is membership in the library's order filtration.
\code{isDifferentialOperator_iff_mem_algebra} then quantifies over $n$.
The theorem \code{challenge_of_twoGeneratorIdentity} places $d$ in the intrinsic
subalgebra, applies the substantive theorem, and coerces its witnesses back to
endomorphisms. Coercion preserves $1$, addition, and composition, and its
injectivity preserves the nonzero hypothesis. No extra mathematical hypothesis
is introduced.

The definitions are restated in
\code{GlobalStafford/Assembly/ChallengeTransport.lean}; \code{GlobalStaffordSolution.lean}
imports the substantive closure and uses the transported theorem. It never
imports \code{GlobalStaffordChallenge.lean}. Comparator checks the Challenge and Solution
statements in separate environments. Its equality check connects the two Lean
statements; this TeX correspondence remains a mathematical reading task.

\section{The two formal phases}\label{sec:phases}
Phase I proves every project-owned implication on concrete carriers. Its endpoint
is \code{GlobalStafford.PhaseI.twoGeneratorIdentity_of_inputs}, taking a value
of \code{GlobalStafford.Assembly.Inputs k A}. The fields specify a finite
principal cover, operator localization and clearance, the whole-chart polynomial
S38 property, chart and global domain structure, and right Ore data.
These are explicit hypotheses, not undeclared axioms.

Phase II constructs every field from Mathlib, AlgebraicAnalysis, and the imported
Stafford38 theorem. The constructor \code{GlobalStafford.inputs} is fed to the
Phase I endpoint by \code{GlobalStafford.universalStatement}. The unconditional
statement therefore has only the field, characteristic, algebra, domain, and
smoothness hypotheses of the theorem above.

\begin{center}
\begin{tabular}{@{}lll@{}}
\toprule
Phase & Packages & Deliverable \\
\midrule
I & WP-1--WP-10 & Conditional descent and chart implications \\
II & WP-11--WP-18 & Concrete producers and universal closure \\
III & WP-19 & Challenge transport and independent replay \\
\bottomrule
\end{tabular}
\end{center}
The record allows only \code{propext}, \code{Classical.choice}, and
\code{Quot.sound} at both terminal endpoints. The single deliberate
\code{sorry} is in the frozen Challenge template; it is absent from the
Solution's import closure. Optional AA-2--AA-4 library extractions concern proofs
already present here and do not reopen Phase I or II.

\section{Weyl S38 over the central parameter field}\label{sec:weyl}
For every characteristic-zero field $K$, every $n\ge0$, and every nonzero
$a\in A_n(K)$, the imported theorem supplies $F,r,s\in A_n(K)$ with
$1=ar+Fas$. The field quantifier is essential: the chart producer invokes it
at $K=k(t)$, not only at the original $k$.

The dependency is \code{itpplasma/stafford38-formal} at
\nolinkurl{e77e176c381ca2d6b20c030f9227f1b031415d2e}, declaration
\code{Stafford38.universalStatement} in \code{Stafford38/FoundationClosure.lean}.
It is a Lean theorem rebuilt against the root project's shared AlgebraicAnalysis
pin, not a literature axiom. Its own proof architecture remains upstream.
The Global Stafford argument uses its universal certificate, not the upstream
fixed-source strengthening.\backmap

\section{Smooth operator structure}\label{sec:structure}
The Phase II producers establish the structural inputs on the actual
algebra $\D_k(A)$. Finite-order operators extend to $A_f$, preserving order;
the extension is injective for nonzero $f$ in the domain $A$, and every
localized operator has a left coefficient denominator. Commutator estimates
then give right denominator clearance.

On an integral affine \'etale chart, coordinate derivations lift uniquely and
commute, finite-order coordinate generation gives normal forms, and the
leading-symbol calculation proves that the operator algebra has no zero
divisors. Weyl right Ore and finite-rank transfer supply chart right Ore;
injective localization and right clearance transfer it back to $\D_k(A)$.
Smoothness supplies a finite principal cover carrying \'etale coordinate maps.
These are the paper's standard structural inputs, discharged in Lean rather
than retained as literature assumptions.

\leanref{GlobalStafford/Assembly/Closure.lean}{GlobalStafford.exists_inputs}

\section{Finite-right-rank ascent}\label{sec:rank}
Let $T\subseteq S$ be domains, with $T$ right Ore, and let $S$ have finite
positive right rank over $T$. If $T$ satisfies S38, then $S$ satisfies S38.
For $0\ne d\in S$, left multiplication by $d$ induces an injective map of the
finite-dimensional right $Q(T)$-space $S\otimes_T Q(T)$. It is surjective,
so clearing a right denominator gives $dy=a$ with $y\in S$, $0\ne a\in T$.
Then
\[
1=ar+Fas=d(yr)+Fd(ys).
\]
No finite generation of $S$ as a $T$-module is required.

Lean expresses finite rank as \code{FiniteRightSpan}: finitely many $s_i\in S$
span every $x\in S$ after multiplying $x$ on the right by some nonzero
denominator from $T$. This is precisely the form produced by chart normal
forms and a finite generic fibre. The localized tensor object is used as a
module; it is not assigned an unjustified ring structure.
\leanref{GlobalStafford/Chart/FiniteRankAscent.lean}{GlobalStafford.Chart.twoGeneratorIdentity_of_finiteRightSpan}

\section{The whole-ring hypothesis on one fixed chart}\label{sec:chart}
Let $C$ be an integral affine algebra \'etale over $B=k[x_1,\ldots,x_n]$.
The coordinate map is injective by flatness and nontriviality. The lifted
coordinate derivations define an injective Weyl action $T=A_n(k)\to S=\D_k(C)$.
Normal forms identify $S$ with $C\otimes_B T$ as a right $T$-module.
Writing $L=\operatorname{Frac}(B)$, the generic fibre $C\otimes_B L$ is a
finite-dimensional field over $L$. Hence $S$ has finite right rank over $T$,
and Section~\ref{sec:rank} proves S38 for $S$.

The same construction works over $k(t)$ on the \emph{same} chart. The formal
scalar extension is
\[
 C_K=k(t)\otimes_k C,\qquad D_K=\D_{k(t)}(C_K).
\]
An injective map $\Theta:S[t]\to D_K$ sends $P$ to its base change and the
central polynomial variable to multiplication by $t$. Every element of $D_K$
is in its image after multiplication by a nonzero scalar polynomial.
Clearing the three witness denominators in a certificate yields
\[
\mathrm{S38}_{\rm poly}(S):\quad
\forall,0\ne E\in S[t],\quad
\exists,0\ne h\in k[t],\ u,w,v\in S[t],\quad h=Eu+wEv.
\]
This is a whole-ring statement, needed for the subsequently constructed
squared input. A certificate for one previously chosen $d$ would not suffice.
The ring permits scalar $t$-denominators only; $t$ is not differentiated.
It is not $\D_k(\operatorname{Frac}(C)(t))$. The \'etale morphism need not be finite.
\leanref{GlobalStafford/Chart/ChartDataOverK.lean}{GlobalStafford.Chart.s38Poly_chart}

\section{Ordered power clearance and contractive powers}\label{sec:order}
Put $R=\D_k(A)$ and $C_f(P)=[P,f]$. If $\ord P\le r$, then
$C_f^{r+1}(P)=0$ and, for $n\ge r$,
\[
Pf^n=\sum_{j=0}^{r}\binom nj f^{n-j}C_f^j(P)
     \in f^{n-r}R_{\le r}.
\]
The identity follows by induction from $Qf=fQ+[Q,f]$.
Thus $f^{-a}P$ has a right denominator cleared by any $f^m$ with $m\ge a+r$.
For $E=f^LP$, $L>r$, induction gives
\[
 E^j\in f^{jL-(j-1)r}R_{\le jr}\qquad(j\ge1).
\]
In the induction step move the fixed-order $P$ past the existing left
coefficient power. Its order is $r$, not the increasing order of $E^j$.
The subsets $f^mR$ are right ideals, not generally two-sided ideals.
\leanref{GlobalStafford/Operators/ContractivePowers.lean}{GlobalStafford.Operators.pow_eq_multiplication_pow_mul}

\section{Squared-annihilator extraction}\label{sec:square}
Fix the current global source $B$ and $0\ne d\in R$. Right Ore supplies
\[
 c=da\ne0,\qquad (Bd)c=db_0.
\]
For $e=chd$, suppose $1=e^2U+We^2V$ in the chart ring. Set
\[
 F=B+Wch,\qquad \alpha=ahd,\qquad\beta=b_0hd.
\]
The required certificate is the exact associative identity
\[
 1=d(\alpha eU-\beta V)+Fd(eV).
\]
Indeed $d\alpha=e$, $d\beta=Bde$, and $Fd=Bd+We$; expansion cancels
$-BdeV+BdeV$ and leaves $e^2U+We^2V$. The square provides the additional
$e$ in the second right cofactor. Every occurrence of $d$ remains in its
specified position.
\leanref{GlobalStafford/Certificate/SquaredAnnihilator.lean}{GlobalStafford.Certificate.certificate_of_square}

\section{Right-moved polynomial conjugation}\label{sec:rho}
The formalization follows paper Remark 4.2. The right binomial identity is
\[
 f^nP=\sum_{j=0}^{r}\binom nj(-1)^j C_f^j(P)f^{n-j},
\]
with terms $j>n$ omitted. In $R_f[t]$ define the finite polynomial
\[
 \rho_t(P)=\sum_j\binom tj(-1)^jC_f^j(P)f^{-j}.
\]
It satisfies $\rho_N(P)=f^NPf^{-N}$ for $N\in\mathbb N$. For the already fixed
Ore data $c,d$, put
\[
 E'(t)=c\rho_t(dc)\rho_{2t}(d),\qquad e_N=cf^Nd.
\]
Then $E'(0)=(cd)^2\ne0$ and $e_N^2=E'(N)f^{2N}$.
A polynomial certificate for $E'$ specializes away from the finitely many
roots of its scalar denominator to a certificate for $e_N^2$.
The source witness is not conjugated in this variant, which simplifies its
uniform order estimate. The full inverse-$\tau_t$ automorphism route in the
paper is not an additional formal endpoint.
\leanref{GlobalStafford/Conjugation/RightMoved.lean}{GlobalStafford.Conjugation.evalNat_squaredInput}

\section{Bounded-order source families}\label{sec:bounded}
Apply $\mathrm{S38}_{\rm poly}(R_f)$ to the fixed nonzero $E'(t)$, obtaining
\[
 h(t)=E'(t)u'(t)+w'(t)E'(t)v'(t),\qquad 0\ne h\in k[t].
\]
At an admissible $N$, take
\[
 U=f^{-2N}u'(N)/h(N),\quad W=w'(N)/h(N),\quad V=f^{-2N}v'(N).
\]
Since the scalar $h(N)$ is central, $1=e_N^2U+We_N^2V$.
Section~\ref{sec:square} gives the successful local source $B+Wcf^N$.
Clear the finitely many coefficients of $w'(t)c$ on the left:
\[
 w'(t)c=f^{-\ell}P(t),\qquad P(t)\in R[t],\qquad \ord P_i\le M.
\]
Both $\ell$ and $M$ are fixed before $N$. The binomial formula makes the source
global for $N\ge\ell+M$:
\[
 F_N=B+f^{N-\ell-M}T_N,\qquad
 T_N=h(N)^{-1}\sum_{j=0}^{M}\binom Nj f^{M-j}C_f^j(P(N)),
 \quad \ord T_N\le M.
\]
Each $P_i$ has order at most $M$, so evaluating the central polynomial does
not increase this bound. Increasing $N$ increases coefficient precision while
the correction order remains bounded. The local cofactors may have growing
orders and pole sizes.
\leanref{GlobalStafford/Descent/BoundedOrderSource.lean}{GlobalStafford.Descent.exists_bounded_order_sources}

\section{Preserving every previously solved chart}\label{sec:protect}
Suppose an old chart $R_g$ has $1=dU_g+BdV_g$. Let the new successful global
source be $F=B+f^LT$, where
\[
 L>M+\ord d+\ord V_g,\qquad \ord T\le M.
\]
Right-clear its new-chart certificate to $f^m=dA_0+FdB_0$ in $R$.
Put $P=TdV_g$, $E=f^LP$, and $r=M+\ord d+\ord V_g$. Then
$dU_g+FdV_g=1+E$. Choose $j$ with $jL-(j-1)r\ge m$; the contractive-powers
lemma gives $(-E)^j=f^mW$ in $R_g$. With the finite sum
$Z=\sum_{i<j}(-E)^i$,
\[
 (1+E)Z=1-(-E)^j,\qquad
 1=d(U_gZ+A_0W)+Fd(V_gZ+B_0W).
\]
This is an identity in the chart ring itself. It uses neither completion nor
an assumption that coefficient ideals are two-sided. At the next chart step,
the actual finite orders of these new cofactors are read before choosing the
next exponent.
\leanref{GlobalStafford/Certificate/OldChartProtection.lean}{GlobalStafford.Certificate.protect_old_chart}

\section{Finite-cover patching}\label{sec:patch}
Take a fixed finite principal cover $f_1,\ldots,f_s$ supplied by smoothness,
with $\mathrm{S38}_{\rm poly}$ on each chart. Start with $B=0$ and no
certificates. At step $i$, retain one global $B$ that works on every earlier
chart. The new chart supplies $M,\ell,H$ from Section~\ref{sec:bounded}.
Only then choose $N$ outside the zeros of $H$ and so large that
\[
 N-\ell-M>M+\ord d+\max_{j<i}\ord V_j.
\]
The new $F_N$ solves chart $i$, and Section~\ref{sec:protect} supplies fresh
certificates on every earlier chart for that same $F_N$.

After exactly $s$ additions, one global $F$ works everywhere on the cover.
Right-clear the local cofactors to
$f_i^{m_i}=dA_i+FdB_i$. These coefficient powers generate the unit ideal in $A$;
choose $a_i\in A$ with $\sum_i f_i^{m_i}a_i=1$. Multiplication on the right and
summation give
\[
 1=d\Bigl(\sum_i A_i a_i\Bigr)+Fd\Bigl(\sum_i B_i a_i\Bigr).
\]
This final B\'ezout argument uses the commutative coefficient algebra only at
the stated point; it does not commute a coefficient through an operator.
\leanref{GlobalStafford/Descent/FiniteCoverPatching.lean}{GlobalStafford.Descent.twoGeneratorIdentity_of_charts}

\section{The terminal assembly}\label{sec:terminal}
Smoothness supplies the finite \'etale principal cover, Section~\ref{sec:chart}
supplies its whole-ring polynomial S38 hypotheses, and
Section~\ref{sec:patch} gives the same-divisor certificate in the original
$\D_k(A)$. These are exactly the ingredients assembled in
\code{GlobalStafford.exists_inputs}. The Phase II theorem then applies the
Phase I implication to those constructed inputs.

The paper also derives torsion-module cyclicity and two-generation of right
ideals. Those consequences are outside this repository's terminal formal
scope. Historical G0--G14 routes, Bjork compression, and the failed uniform
cofactor-lifting approaches are not dependencies of this proof.
\leanref{GlobalStafford/Assembly/Closure.lean}{GlobalStafford.universalStatement}

\section{Verification and reproducible dependencies}\label{sec:verification}
The repository report \code{docs/verification-results.json} records the
isolated Phase I/full replay and Comparator acceptance on 9 September 2026.
Both NanoDa and Lean accepted the solution. The report SHA-256 is
\nolinkurl{08b8dcd152806c5033316100a624d8c61a67f9e069c18ac649f3c6612779db67}.
A second comparison using Palomar's canonical Challenge compiler and protected
alias also passed. Its outer systemd wrapper was unavailable locally;
Landrun and the mathematical checks were retained. Hosted registry acceptance
and human expert review are separate.

\begin{description}
\item[Lean] \code{leanprover/lean4:v4.33.0}
\item[Mathlib] \nolinkurl{db584cd6d46c92f209a44c0f1c829460d327499d}
\item[AlgebraicAnalysis] \nolinkurl{4aae47967f6ba02ffe2f639ab06564c9a9d1ecc8}
\item[Stafford38 formal] \nolinkurl{e77e176c381ca2d6b20c030f9227f1b031415d2e}
\end{description}
Lake uses the root manifest's AlgebraicAnalysis revision for both projects.
The verifier builds \code{Stafford38.FoundationClosure} and
\code{Stafford38.LocalizedDifferentialCorollaries} against that shared pin.
A newer upstream release is not substituted without a dependency work package.

\begin{Verbatim}[fontsize=\small,frame=single]
lake exe cache get
lake build
scripts/verify.sh --phase-i
scripts/verify.sh
scripts/bootstrap-palomar-tools.sh
scripts/verify-palomar.sh
\end{Verbatim}
Run the verifier modes sequentially and retain their logs before the next
mode overwrites \code{.lake/verification}. The paper's 984 exact rational
regression assertions check ordered identities on finite examples; they do
not replace universal Lean proofs or human review.

\section{Release companion and Zenodo description}\label{sec:release}
This companion follows the Stafford38 Challenge dossier and proof-map
supplement. Stafford38 \code{v1.1.0} includes the stronger exact-source comparison and corrected proof architecture. Both its Comparator configurations passed NanoDa and Lean kernel checking on the controller host. Its concept DOI is
\href{https://doi.org/10.5281/zenodo.22390721}{10.5281/zenodo.22390721}.

For Global Stafford, a source release can include this dossier's TeX,
clickable map, PDF and checksums, together with the verification report and
pinned software source. The prepared Zenodo description is:
\begin{quote}
Lean 4 formalization of the identity $1=dR+FdS$ for every nonzero finite-order
$k$-linear differential operator $d$ on a smooth integral affine algebra over
any characteristic-zero field. The proof extends the formalized Weyl case by
fixed \'etale charts, bounded-order source corrections, preservation of earlier
chart certificates, and finite-cover descent. Phase II constructs all
structural inputs from Mathlib and the pinned AlgebraicAnalysis and
Stafford38 developments. The package includes the Challenge/Solution pair,
verification evidence, and a companion document with a clickable proof map,
the statement correspondence, and the main proof mechanisms.
\end{quote}
The formal repository is public under Apache-2.0. Release \code{v1.0.5} packages the corrected provenance, dated literature audit and registry citation metadata. The mathematical replay is unchanged. Human expert review is separate.

\appendix
\clearpage
\section{The actual Challenge file}\label{sec:challenge}
This listing is read directly from the repository file at build time.
Its deliberate final \code{sorry} is the statement placeholder for comparison.
\VerbatimInput[fontsize=\scriptsize,breaklines=true,breakanywhere=true,
 frame=lines,numbers=left,numbersep=5pt]{../../GlobalStaffordChallenge.lean}

\clearpage
\section{The actual Solution file and Comparator configuration}\label{sec:solution}
The Solution reaches the transported theorem through the substantive closure.
The transport's recursive equivalences are explained in
Section~\ref{sec:comparison}.
\VerbatimInput[fontsize=\small,breaklines=true,breakanywhere=true,
 frame=lines,numbers=left,numbersep=5pt]{../../GlobalStaffordSolution.lean}
\VerbatimInput[fontsize=\small,breaklines=true,breakanywhere=true,
 frame=lines]{../../comparator.json}

\section{Literature and proof sources}\label{sec:literature}
Documentation release v1.0.6, 9 September 2026. Mathematical declarations and dependency pins are unchanged.

The Weyl node uses the proved Stafford38 theorem at the manifest pin;
\cite{Sta78} supplies the conjecture. Smooth operator structure has the
background \cite{StacksD,StacksS,Bav10}. Finite-rank ascent, squared-annihilator
extraction, bounded-order sources and chart protection are project proofs.
The three-generator result \cite{CH88} and survey \cite{Bel26} supply prior-art
context, not assumptions of the descent.

Full source roles and clickable Lean entry points are indexed at \url{https://github.com/itpplasma/global-stafford-formal/blob/main/docs/literature.md}.

\begin{thebibliography}{StacksS}
\bibitem[Sta78]{Sta78} J. T. Stafford, \emph{Module Structure of Weyl Algebras}, Journal of the London Mathematical Society (2) 18 (1978), 429–442. \url{https://doi.org/10.1112/jlms/s2-18.3.429}.
\bibitem[CH88]{CH88} S. C. Coutinho and M. P. Holland, \emph{Module Structure of Rings of Differential Operators}, Proceedings of the London Mathematical Society (3) 57 (1988), 417–432. \url{https://doi.org/10.1112/plms/s3-57.3.417}.
\bibitem[Bel26]{Bel26} Gwyn Bellamy, \emph{Module structure of Weyl algebras}, Journal of the London Mathematical Society 113 (2026), no. 1, e70373. \url{https://doi.org/10.1112/jlms.70373}.
\bibitem[StacksD]{StacksD} The Stacks Project Authors, \emph{Finite order differential operators}, The Stacks Project, Section 10.133, tag 09CH (accessed 2026-09-09). \url{https://stacks.math.columbia.edu/tag/09CH}.
\bibitem[StacksS]{StacksS} The Stacks Project Authors, \emph{Smooth ring maps}, The Stacks Project, Section 10.137, tag 00T1 (accessed 2026-09-09). \url{https://stacks.math.columbia.edu/tag/00T1}.
\bibitem[Bav10]{Bav10} V. V. Bavula, \emph{Generators and Defining Relations for the Ring of Differential Operators on a Smooth Affine Algebraic Variety}, Algebras and Representation Theory 13 (2010), 159–187. \url{https://doi.org/10.1007/s10468-008-9112-7}.
\end{thebibliography}
\end{document}
